\section{Discussion}
\label{sec:discussion}

\subsection{Implications for AI System Evaluation}

The \ARES{} framework provides a principled basis for evaluating AI agent
reliability \textit{before deployment}. Because four of the five mismatch
dimensions ($\delta_S$, $\delta_K$, $\delta_C$, $\delta_T$) are computable
from the agent's declared profile and the task's requirements alone, ARS
can serve as a pre-execution deployment gate: agents scoring below $\theta^*$
are automatically escalated to human analysts rather than executing
autonomously. This approach is consistent with the NIST AI Risk Management
Framework~\cite{nist_ai_rmf}, which recommends continuous monitoring of AI
system performance against operational requirements.

The task-class-specific weight vectors (Table~\ref{tab:weights_isaia})
provide additional operational value: they reveal which capability
dimensions matter most for each threat type. A SOC deploying an agent for
ransomware detection should prioritise knowledge coverage ($w_K = 0.35$),
while an agent assigned to DDoS mitigation should be evaluated primarily
on skill and capability access ($w_S + w_C = 0.60$). This dimension-specific
guidance enables targeted agent improvement rather than generic capability
enhancement.

\subsection{Trustworthy Deployment of AI Agents}

The Epistemic Gap metric supports runtime monitoring of hallucination
risk: by comparing an agent's stated confidence against its known
knowledge coverage, organisations can flag high-risk outputs for human
review before they trigger automated actions. The temporal drift model
($\delta_T$) provides a quantitative basis for scheduling knowledge
updates -- e.g., for high-volatility domains like ransomware
($\lambda = 0.005$), ARS degrades below $\theta^*$ in approximately
$120$ days, suggesting a quarterly refresh cycle. The cascade
reliability bound (Theorem~\ref{thm:cascade}) implies that decomposing
a task into a tightly-coupled pipeline cannot improve reliability beyond
the weakest agent, so single-agent designs should be preferred for
critical tasks unless per-agent reliability is independently verified.

\subsection{Limitations}

\textbf{LLM specificity.} Experiments use Mistral-7B-Instruct; absolute
thresholds will require recalibration for larger models, though directional
findings should hold since the SKC model is architecture-agnostic.
\textbf{Linear formulation.} The additive ARS is chosen for
interpretability and decomposability; nonlinear alternatives may improve
predictive power at the cost of transparency.
\textbf{Synthetic data.} The dataset is generated from ATT\&CK and NVD;
real-world SOC incidents introduce additional noise and ambiguity.
\textbf{Cascade model.} The formulation assumes acyclic pipelines and
independent mismatch factors; cyclic architectures and correlated
dimensions require further modelling.
\textbf{Controlled benchmark.} Because the benchmark uses controlled
mismatch injection, real SOC deployments may exhibit noisier and
less separable failure patterns.
